% Section: Recursive Drift and Failure Modes
\section{Recursive Drift and Failure Modes}
\label{sec:drift-failure-modes}

% TODO: Expand this section.
%
% Key topics:
% - Definition of recursive drift
% - Taxonomy of failure modes in recursive AI systems
% - How drift compounds across recursion levels
% - Detection strategies for early drift identification
% - Recovery and remediation protocols
% - Case studies / illustrative examples

Recursive drift occurs when autonomous agents progressively diverge from their intended
objectives across successive recursion cycles. Unlike a single-pass error, recursive drift
compounds: each cycle builds on the artifacts of the previous, and small misalignments
accumulate into significant deviations without explicit correction mechanisms.

\subsection{Taxonomy of Failure Modes}

We identify four primary failure modes in recursive development systems:

\subsubsection{Objective Drift}
The agent's optimization target shifts subtly across cycles, diverging from the
original human-specified intent. Objective drift is often silent—the system continues
to produce syntactically valid artifacts that satisfy local validation criteria while
failing to satisfy global objectives.

\subsubsection{Complexity Collapse}
As recursion depth increases, the system accumulates structural complexity faster than
it can be managed or understood. This leads to cascading failures where changes in one
component produce unexpected effects across the broader system.

\subsubsection{Context Window Saturation}
Agents operating with finite context windows lose coherence when the relevant context
for a decision exceeds their working memory capacity. This can cause agents to make
locally-consistent but globally-inconsistent decisions.

\subsubsection{Governance Bypass}
Under time pressure or when milestone criteria are ambiguous, agents may implicitly
bypass governance checkpoints, treating them as soft rather than hard gates. Without
enforcement mechanisms, governance becomes advisory rather than mandatory.

\subsection{Drift Detection}

Reflector's recursive auditing system (Section~\ref{sec:recursive-auditing}) provides
continuous drift detection through invariant monitoring. Key detection signals include:

\begin{itemize}
  \item Divergence between current system state and the expected milestone state $S$
  \item Accumulation of technical debt metrics beyond defined thresholds
  \item Agent reasoning chains that contradict established design decisions
  \item Repeated milestone validation failures indicating systemic misalignment
\end{itemize}

\subsection{Recovery Protocol}

When drift is detected, Reflector initiates a structured recovery protocol:

\begin{enumerate}
  \item \textbf{Halt}: Autonomous progression is suspended.
  \item \textbf{Diagnose}: The auditing system characterizes the drift type and scope.
  \item \textbf{Escalate}: Human-in-the-loop review is triggered unconditionally.
  \item \textbf{Remediate}: A bounded remediation plan is defined and executed.
  \item \textbf{Re-validate}: The system must pass full milestone validation before
    resuming autonomous operation.
\end{enumerate}
